\section{Mixture-of-Experts and Adapter Routing}
\label{sec:moe_routing}

This chapter examines mixture-of-experts architectures and adapter routing
mechanisms, asking how existing routing approaches select among a known set of
expert modules and where their assumptions break down in the context of a
decentralised, dynamically evolving adapter pool. The routing logic reviewed
here is technically relevant to the research direction examined in this SLR;
the gap is not in the routing mechanism itself but in the prerequisite catalogue
that all reviewed routing systems silently assume.

\subsection{Mixture-of-Experts as a Scaling Principle}
\label{sec:moe_routing:foundations}

Mixture-of-Experts (MoE) architectures \citep{ShazeerMoE2017} replace dense computation with \textit{conditional}
computation: a routing function selects a subset of expert modules to activate for each
input, allowing the total model capacity to grow without a proportional increase in
inference cost. The MoE principle is directly relevant to adapter serving: if each
PEFT adapter is treated as one expert, a routing mechanism can select the appropriate
adapter for each incoming request without loading all adapters simultaneously, which is an
efficiency property that scales to large adapter libraries.

\citet{Fedus2022} introduced Switch Transformers, which replace the feedforward
sub-layer in each transformer block with a sparse MoE layer whose \textit{router}
assigns each token to exactly one expert (Top-1 routing). Top-1 routing eliminates the
token-drop problem of earlier Top-$k$ schemes by introducing a configurable \textit{expert
capacity factor} as a fraction of the per-expert buffer, reducing token dropping to near zero. Switch Transformers achieved
a 7$\times$ pre-training speedup over dense T5 at equivalent parameter count, scaling
to trillion-parameter sparse models. A load-balancing auxiliary loss prevents routing
collapse in which a small subset of experts receives all tokens. However, Switch
Transformers rely on a fixed set of known experts co-located in a single process; they do
not address discovering new experts from remote peers or routing across a P2P network
deploy. This gap is examined in the synthesis and gap analysis
(Section~\ref{sec:synthesis_gap}).

\citet{Zhou2022} proposed Expert Choice routing, which inverts the routing perspective:
rather than each token selecting its top expert, each expert selects its top-$k$ tokens
from the batch. This guarantees perfect load balancing by construction and eliminates
token overflow, at the cost of requiring global batch visibility at the router.
Expert Choice consistently outperforms token-choice routing on downstream NLP benchmarks
at matched capacity.

\citet{RajbhandariDeepSpeed2022} introduced DeepSpeed-MoE, an end-to-end training and
inference framework for MoE models on heterogeneous hardware, addressing the engineering
challenges of all-to-all communication required for expert dispatch across multiple GPUs.
DeepSpeed-MoE demonstrated that MoE inference latency is dominated by expert dispatch
communication, not computation, and proposed hierarchical expert placement strategies
to minimise cross-node communication.

\subsection{Towards Distributed and Crowdsourced MoE}
\label{sec:moe_routing:distributed}

The MoE literature has begun to address the challenge of distributing experts across
geographically separated nodes rather than co-locating them within a single cluster.
\citet{RyabininCrowdsourced2020} proposed a crowdsourced MoE architecture in which
expert modules are hosted by volunteer participants over the internet, introducing
the problems of peer churn, heterogeneous latency, and Byzantine expert behaviour.
The authors demonstrated that crowdsourced MoE can serve useful inference despite
dropout rates of up to 30\% of experts per batch, by over-provisioning the expert
pool and routing around failed nodes.

\citet{ZadouriPushMoE2024} proposed pushing MoE to its practical limits through
adaptive expert specialisation under communication constraints, demonstrating that
MoE routing quality degrades gracefully when all-to-all communication is limited
to $O(\log N)$ messages per batch --- a communication budget consistent with DHT-based
lookup.

While the MoE literature reviewed above operates within a single training or inference
run, the Petals system of \citet{Borzunov2023} demonstrated routing at the
\textit{node level}: client requests are dispatched to whichever volunteer nodes
hold the required transformer blocks, with a DHT-based routing table mapping block
indices to peer addresses. This node-level routing is structurally analogous to
adapter-level routing in a P2P exchange architecture, the DHT maps requested adapters to
the peers that hold them, and the client fetches from the closest or fastest provider.

\subsection{LoRA-Level MoE: Routing Among Adapter Experts}
\label{sec:moe_routing:lora_moe}

A rapidly growing body of work applies the MoE routing principle directly to PEFT
adapters, treating each LoRA matrix or bottleneck adapter as one expert within a
sparse routing framework. This shifts the routing granularity from the layer level
(Switch Transformers) to the adapter level, and is directly applicable to the problem
of selecting adapters from a library in a multi-task serving system.

These methods differ primarily in routing granularity and gating mechanism.
At the finest granularity, MOELoRA \citep{LiuMOELoRA2023} attaches separate LoRA experts
per attention head and routes each token to its top-$k$ experts at the head level,
allowing different subspaces of the key--value projection to specialise for different tasks.
The remaining methods operate at the layer level but vary in how routing weights are
computed: MixLoRA \citep{LiMixLoRA2024} uses continuous mixing weights produced by a
differentiable gating function with contrastive regularisation to prevent expert collapse,
while MiLoRA \citep{ZhangMiLoRA2024} uses Top-1 sparse gating with a diversity penalty
that encourages each expert to capture a distinct region of the task distribution.
MoELoRA \citep{LuoMoELoRA2024} and Mixture of LoRA Experts \citep{WuMixtureLoRA2024}
independently demonstrated that routing among multiple LoRA experts within a single
fine-tuning run outperforms standard LoRA and full fine-tuning on multi-task benchmarks,
using different gating mechanisms to achieve this result.

A common finding across all reviewed LoRA-MoE methods is that multi-expert routing
consistently outperforms single LoRA at matched parameter count. Beyond raw accuracy,
two specialised concerns have been addressed: LoRAMoE \citep{DouLoRAMoE2024} targets the
world knowledge preservation problem in instruction-tuned models by routing
world-knowledge queries to a subset of frozen LoRA experts, and MoLoRAs
\citep{FengMoLoRAs2024} demonstrates that conditioning the gating function on both task
identity and input embedding consistently outperforms input-only routing across
domain-diverse benchmarks. MixLoRA's sparse top-$k$ routing is additionally noted as
more stable than uniform averaging under small batch sizes, making it suitable for
low-throughput deployment scenarios.

\subsection{Mixture-of-Domain Adapters and Broader Routing Policies}
\label{sec:moe_routing:domain}

Beyond LoRA-specific architectures, a complementary line of work applies MoE-style
routing to domain adapter modules. \citet{DiaomMoDE2023} proposed Mixture-of-Domain
Experts (MoDE), which maintains separate adapter modules per data domain and uses a
soft gating network to interpolate domain-adapter outputs at inference time. MoDE
enables a single model to serve multiple domains without per-domain loading, achieving
accuracy within 1--2\% of per-domain fine-tuning.

\citet{WangAdaMix2022} proposed AdaMix, a general mixture-of-adaptations framework
that routes between multiple heterogeneous PEFT modules (combining bottleneck adapters
and LoRA matrices within the same layer) using a sparse learned gating network. AdaMix
outperforms any single PEFT method on GLUE and SuperGLUE at matched parameter budgets,
demonstrating that heterogeneous PEFT composition via routing is superior to any
homogeneous single method.

\citet{Muqeeth2024} investigated soft merging of task-specific adapter experts, showing
that interpolating adapter weights according to input-conditioned routing weights
outperforms both hard routing and uniform averaging, with a particularly strong
advantage on distribution-shifted test sets.

\citet{LiuWhenMOE2024} provided a comprehensive analysis of when MoE routing benefits
LLM fine-tuning, identifying task diversity and data heterogeneity as the primary
drivers of MoE gain over single-adapter baselines, and routing collapse as the primary
failure mode when task distribution is narrow.

\subsection{Sparse Adapters and Efficient Routing}
\label{sec:moe_routing:sparse}

\citet{LiZhouSparser2025} proposed SMOA, a sparse mixture-of-adapters architecture that
uses a cross-layer shared adapter pool and a global router to select a sparse subset of
experts per layer, reducing active adapters by over 85\% while preserving task accuracy.
The architecture introduces a global pool of adapters shared across layers, and a
curriculum learning strategy that transitions adapters from per-layer specialisation to
cross-layer generalisation. Routing is performed by a global router that selects the
top-activated adapters per layer using a scoring function over all adapters in the
pool, with the condition that at most $n_l$ adapters are active per layer. SMOA also
incorporates expert-redundancy regularisation that steers adapters toward learning
specialised residual knowledge complementary to the backbone, using the backbone's
weight in the routing decision as a regulariser. However, the routing is designed for
centralised multi-task settings with a known, fixed adapter pool; it does not address
discovering new experts from remote peers or routing in a network where both the expert
pool and routing decisions must be learned in a distributed manner.


\subsection{Where MoE Routing Stops Short}
\label{sec:moe_routing:implications}


All routing solutions reviewed in this chapter share one silent prerequisite: the expert pool must be fully known and enumerable before any routing decision can be made. In co-located, fixed-topology systems this prerequisite is trivial — it is simply another configuration parameter, so obvious it rarely needs stating. In a P2P network, where peer presence is dynamic, depends on factors outside any single node's control, and where adapters are contributed without central coordination, that same trivial parameter becomes an open and unsolved problem.

The routing algorithms themselves — Top-k gating, Expert Choice, differentiable mixing — would function correctly in a P2P setting once supplied with a complete catalogue over which to route (even crowdsourced MoE (\citet{RyabininCrowdsourced2020}) relies on a pre-registered volunteer pool). In this sense the gap is not in the routing logic or mechanism. The gap is in the catalogue itself: how to construct it dynamically from a distributed, self-evolving, and unpredictable set of peers whose adapter holdings change without notice.

All reviewed routing methods assume a fully enumerated expert pool. In a P2P setting, this assumption does not hold: no single node would know the full set of available adapters. This raises the question of whether compact, data-free capability representations — derived from adapter weights alone — could supply the missing catalogue information that existing routing mechanisms require. Whether such representations carry sufficient task-level signal for similarity-based retrieval remains an open question, examined in Section~\ref{sec:synthesis_gap}.

